% Policy-Based Data Structures (PBDS)

An extension structure in GCC that provides a high-performance Ordered Set, supporting index-based lookups (\lstinline|find_by_order|) and order queries (\lstinline|order_of_key|) in $O(\log N)$.

\begin{lstlisting}
#include <ext/pb_ds/assoc_container.hpp>
#include <ext/pb_ds/tree_policy.hpp>
using namespace __gnu_pbds;

// Definition for Ordered Set
template <typename T>
using ordered_set = tree<T, null_type, less<T>, rb_tree_tag, tree_order_statistics_node_update>;

// Definition for Ordered Multiset (allows duplicate keys)
template <typename T>
using ordered_multiset = tree<T, null_type, less_equal<T>, rb_tree_tag, tree_order_statistics_node_update>;
\end{lstlisting}

Basic operations and functions usage pattern:
\begin{lstlisting}
ordered_set<int> st;

// 1. Insert elements normally
st.insert(1); st.insert(4); st.insert(8);

// 2. find_by_order(k): Returns an iterator to the k-th smallest element (0-indexed)
auto it = st.find_by_order(1); // Points to '4'
int val = *it;

// 3. order_of_key(x): Returns the number of elements strictly smaller than x
int cnt = st.order_of_key(5); // Returns 2 (elements 1 and 4 are smaller)

// 4. Erase element in ordered_multiset (MUST use iterator to avoid removing all duplicates)
ordered_multiset<int> mst;
mst.insert(4); mst.insert(4);
mst.erase(mst.find_by_order(mst.order_of_key(4))); // Safely erases only one '4'
\end{lstlisting}
